\paragrafosz{Εισαγωγή: οι φυσικοί αριθμοί και οι πράξεις τους}{Α.1.0, «Εισαγωγή»}

\begin{anagke}
Η κυρία Αλεξάνδρα γράφει στον πίνακα τρία ερωτήματα, χωρίς να πει τίποτα άλλο. «Ποιος αριθμός έρχεται μετά το 999;» «Πόσο κάνει $4+3\cdot 5$;» «Τι σημαίνει το 5 στον αριθμό 5.021;» Η Ζωή απαντά αμέσως στο πρώτο. Ο Νίκος βιάζεται και λέει $35$ για το δεύτερο, ενώ η Φανή γράφει $19$. Για το τρίτο ο Ηλίας σηκώνει το χέρι: «Πέντε χιλιάδες, αφού είναι στην πρώτη θέση».

Όλοι αυτοί οι αριθμοί και οι πράξεις τούς είναι γνωστοί από το Δημοτικό. Τώρα όμως θέλουμε να τα δούμε πιο προσεκτικά, γιατί σε όλο το βιβλίο θα στηριζόμαστε σε αυτά: κάθε φορά που θα έχουμε αμφιβολία, θα γυρίζουμε σε αυτή τη σελίδα.

Σε αυτή την παράγραφο θυμόμαστε τους φυσικούς αριθμούς, την αξία των ψηφίων τους, τη διάταξή τους, τις ιδιότητες των πράξεων και τη σειρά με την οποία γίνονται οι πράξεις.
\end{anagke}

\begin{lexi}{φυσικός, άρτιος, περιττός, προσθετέος, παράγοντας}{φυσικοί αριθμοί}
Οι αριθμοί λέγονται \lexiword{φυσικοί} γιατί προέκυψαν «φυσικά» από το μέτρημα: ένα, δύο, τρία. Ο \lexiword{άρτιος} (από το «αρτίως», δηλαδή «ακριβώς») χωρίζεται δύο-δύο χωρίς να περισσεύει τίποτα. Ο \lexiword{περιττός} είναι αυτός που «περισσεύει» όταν τον μοιράσεις σε ζευγάρια. Ο \lexiword{προσθετέος} είναι αυτός που «πρόκειται να προστεθεί», ενώ οι \lexiword{παράγοντες} είναι όσοι «παράγουν», δηλαδή φτιάχνουν, το γινόμενο.
\end{lexi}

\begin{metafrasi}
Ο αριθμός μετά το 9 είναι ο επόμενός του. & $9+1=10$\\
Πέντε χιλιάδες, δύο δεκάδες και μία μονάδα. & $5.021$\\
Ο τρεις είναι μικρότερος από τον πέντε. & $3<5$\\
Η σειρά που προσθέτω δεν αλλάζει το αποτέλεσμα. & $\alpha+\beta=\beta+\alpha$\\
Πρώτα το πολλαπλασιάζω, μετά το προσθέτω. & $4+3\cdot 5=4+15=19$\\
\end{metafrasi}

\tmima{Η θεωρία}

\begin{orismos}{φυσικοί αριθμοί}
Οι αριθμοί $0,1,2,3,4,5,6,7,8,9,10,11,12,\dots$ ονομάζονται \textbf{φυσικοί αριθμοί}. Κάθε φυσικός έχει έναν προηγούμενο και έναν επόμενο, εκτός από το $0$, που έχει μόνο επόμενο.
\end{orismos}

Οι φυσικοί χωρίζονται σε δύο ομάδες. \textbf{Άρτιοι} (ή ζυγοί) είναι αυτοί που διαιρούνται με το $2$: $0,2,4,6,8,10,\dots$ \textbf{Περιττοί} (ή μονοί) είναι οι άλλοι: $1,3,5,7,9,11,\dots$ Οι άρτιοι τελειώνουν σε $0,2,4,6$ ή $8$ και οι περιττοί σε $1,3,5,7$ ή $9$.

\tmimasmall{Το δεκαδικό σύστημα και η αξία των ψηφίων}

Γράφουμε όλους τους φυσικούς με τα δέκα ψηφία $0,1,2,\dots,9$. Λέμε ότι το σύστημά μας είναι το \textbf{δεκαδικό} και έχει \textbf{βάση το δέκα}. Η αξία κάθε ψηφίου εξαρτάται από τη θέση του, δηλαδή από την τάξη στην οποία βρίσκεται. Στον αριθμό $1.253.408$ έχουμε από αριστερά: 1 μονάδα εκατομμυρίων, 2 εκατοντάδες χιλιάδων, 5 δεκάδες χιλιάδων, 3 μονάδες χιλιάδων, 4 εκατοντάδες, 0 δεκάδες και 8 μονάδες.

\begin{center}
\renewcommand{\arraystretch}{1.5}\setlength{\tabcolsep}{5pt}
\small
\begin{tabular}{|c|c|c|c|c|c|c|}\hline
\scriptsize εκατομμύρια & \scriptsize εκατ.\ χιλιάδων & \scriptsize δεκ.\ χιλιάδων & \scriptsize χιλιάδες & \scriptsize εκατοντάδες & \scriptsize δεκάδες & \scriptsize μονάδες\\\hline
1 & 2 & 5 & 3 & 4 & 0 & 8\\\hline
\end{tabular}
\end{center}

\tmimasmall{Διάταξη και σύγκριση}

Οι φυσικοί μπορούν να διαταχθούν από τον μικρότερο προς τον μεγαλύτερο: $0<1<2<3<4<\dots$ Τους τοποθετούμε πάνω σε μια ευθεία, την αριθμογραμμή: επιλέγουμε ένα σημείο για το $0$, δεξιά του ένα δεύτερο σημείο για το $1$ και με την ίδια απόσταση βρίσκουμε τα σημεία για τα $2,3,4,\dots$ Όσο πηγαίνουμε προς τα δεξιά, οι αριθμοί μεγαλώνουν.

\begin{center}
\akgr[0.62]{0}{15}{3/,8/,12/}
\end{center}

Για να συγκρίνουμε δύο αριθμούς, βλέπουμε πρώτα πόσα ψηφία έχουν (ο πιο μακρύς είναι μεγαλύτερος) και, αν έχουν τα ίδια, συγκρίνουμε τα ψηφία από αριστερά προς τα δεξιά.

\tmimasmall{Οι πράξεις και οι ιδιότητές τους}

Στην πρόσθεση $\alpha+\beta=\gamma$ οι $\alpha,\beta$ λέγονται \textbf{προσθετέοι} και ο $\gamma$ \textbf{άθροισμα}. Στην αφαίρεση $\mu-\alpha=\delta$ ο $\mu$ λέγεται \textbf{μειωτέος}, ο $\alpha$ \textbf{αφαιρετέος} και ο $\delta$ \textbf{διαφορά}. Στους φυσικούς η αφαίρεση γίνεται μόνο όταν ο αφαιρετέος είναι μικρότερος ή ίσος από τον μειωτέο. Στον πολλαπλασιασμό $\alpha\cdot\beta=\gamma$ οι $\alpha,\beta$ λέγονται \textbf{παράγοντες} και ο $\gamma$ \textbf{γινόμενο}. Το σύμβολο του πολλαπλασιασμού συχνά παραλείπεται: γράφουμε $2\alpha$ αντί για $2\cdot\alpha$ και $3(1+4)$ αντί για $3\cdot(1+4)$.

\begin{idiotita}{πρόσθεση}
$\alpha+\beta=\beta+\alpha$ (\emph{αντιμεταθετική}) \qquad $\alpha+(\beta+\gamma)=(\alpha+\beta)+\gamma$ (\emph{προσεταιριστική}) \qquad $\alpha+0=\alpha$ (το $0$ είναι το \emph{ουδέτερο στοιχείο} της πρόσθεσης).
\end{idiotita}

\begin{idiotita}{πολλαπλασιασμός}
$\alpha\cdot\beta=\beta\cdot\alpha$ (\emph{αντιμεταθετική}) \qquad $\alpha\cdot(\beta\cdot\gamma)=(\alpha\cdot\beta)\cdot\gamma$ (\emph{προσεταιριστική}) \qquad $\alpha\cdot 1=\alpha$ (το $1$ είναι το \emph{ουδέτερο στοιχείο} του πολλαπλασιασμού) \qquad $\alpha\cdot 0=0$ (το $0$ είναι το \emph{απορροφητικό στοιχείο}).\\
\textbf{Επιμεριστική ιδιότητα:} $\alpha\cdot(\beta+\gamma)=\alpha\cdot\beta+\alpha\cdot\gamma$ και $\alpha\cdot(\beta-\gamma)=\alpha\cdot\beta-\alpha\cdot\gamma$.
\end{idiotita}

\tmimasmall{Αριθμητικές παραστάσεις και προτεραιότητα των πράξεων}

Μια παράσταση που περιέχει πράξεις με αριθμούς λέγεται \textbf{αριθμητική παράσταση}, όπως η $2\cdot 3+7-1$ ή η $(15-6)\cdot(2^3-9\cdot 0)$. Το αποτέλεσμα όλων των πράξεων λέγεται \textbf{τιμή της παράστασης}.

\begin{kanonas}{προτεραιότητα των πράξεων}
\begin{enumerate}[leftmargin=1.6em,itemsep=0.2ex,topsep=2pt]
\item Υπολογίζουμε πρώτα τις \textbf{δυνάμεις}.
\item Μετά τους \textbf{πολλαπλασιασμούς και τις διαιρέσεις}, με τη σειρά που εμφανίζονται, από αριστερά προς τα δεξιά.
\item Τέλος τις \textbf{προσθέσεις και τις αφαιρέσεις}, πάλι από αριστερά προς τα δεξιά.
\end{enumerate}
Αν υπάρχουν παρενθέσεις, εκτελούμε πρώτα τις πράξεις μέσα σε αυτές, με την παραπάνω σειρά.
\end{kanonas}

(Τις δυνάμεις θα τις γνωρίσουμε αναλυτικά στην επόμενη παράγραφο.)

\begin{prosoxi}
\begin{itemize}[leftmargin=1.2em,itemsep=0.3ex,topsep=2pt]
\item Το $4+3\cdot 5$ δεν κάνει $35$ αλλά $19$: ο πολλαπλασιασμός προηγείται.
\item Στην αφαίρεση και τη διαίρεση δεν ισχύουν οι ιδιότητες των άλλων δύο πράξεων. Για παράδειγμα $(8-3)-1=4$, ενώ $8-(3-1)=6$.
\end{itemize}
\end{prosoxi}

\tmima{Λυμένα παραδείγματα}

\paradeigma{Αξία των ψηφίων}
\tixero Ο αριθμός $4.063.518$.\par
\tithelo Την αξία κάθε ψηφίου και τον επόμενο και τον προηγούμενο του αριθμού.\par
\possindeo Διαβάζω από δεξιά: 8 μονάδες, 1 δεκάδα, 5 εκατοντάδες, 3 χιλιάδες, 6 δεκάδες χιλιάδων, 0 εκατοντάδες χιλιάδων, 4 εκατομμύρια. Ο επόμενος προκύπτει αν προσθέσω 1, ο προηγούμενος αν αφαιρέσω 1.\par
\apantisi Επόμενος $4.063.519$, προηγούμενος $4.063.517$.\par
\elegxos Το $0$ στις εκατοντάδες χιλιάδων δεν λείπει: κρατά τη θέση του, ώστε το $4$ να μείνει στα εκατομμύρια.

\paradeigma{Η επιμεριστική ιδιότητα στον νου}
\tixero Θέλω το γινόμενο $7\cdot 102$ χωρίς κομπιουτεράκι.\par
\tithelo Έναν γρήγορο υπολογισμό.\par
\possindeo Γράφω $102=100+2$ και εφαρμόζω την επιμεριστική: $7\cdot(100+2)=7\cdot 100+7\cdot 2=700+14=714$.\par
\apantisi $714$.\par
\elegxos Ο ίδιος υπολογισμός με $102\cdot 7=714$ (αντιμεταθετική).

\paradeigma{Η παράσταση με πολλές πράξεις}
\tixero Η παράσταση $40-(3+5)\cdot 4+12:6$.\par
\tithelo Την τιμή της.\par
\possindeo Πρώτα η παρένθεση: $3+5=8$. Μετά οι πολλαπλασιασμοί και οι διαιρέσεις: $8\cdot 4=32$ και $12:6=2$. Έχω $40-32+2$. Τέλος από αριστερά προς τα δεξιά: $40-32=8$ και $8+2=10$.\par
\apantisi $10$.\par
\elegxos Αν είχα προσθέσει πρώτα το $2$ στο $32$ θα έβγαζα $40-34=6$. Η σωστή σειρά είναι από αριστερά προς τα δεξιά.

\tmima{Ασκήσεις}

\zoni{Το κατάλαβα}{θυμάμαι και ονομάζω}

\ex Γράψε τον προηγούμενο και τον επόμενο των αριθμών $0$, $99$, $1.000$ και $25.099$.

\ex Στον αριθμό $4.635.182$ γράψε ποια τάξη έχει κάθε ψηφίο και ποια είναι η αξία του $6$ και του $3$.

\ex Συμπλήρωσε με $<$, $>$ ή $=$: α) $3.099\ \dots\ 3.101$ \quad β) $45.670\ \dots\ 45.607$ \quad γ) $1.000.000\ \dots\ 999.999$ \quad δ) $7.250\ \dots\ 7.250$.

\ex Ποια ιδιότητα εφαρμόζεται κάθε φορά; α) $7+12=12+7$ \quad β) $(4+5)+6=4+(5+6)$ \quad γ) $9\cdot 1=9$ \quad δ) $13\cdot 0=0$ \quad ε) $6\cdot(10+3)=6\cdot 10+6\cdot 3$.

\zoni{Το χρησιμοποιώ}{πράξεις και παραστάσεις}

\ex \textbf{Νοερός υπολογισμός.} Υπολόγισε χωρίς κομπιουτεράκι με την επιμεριστική ιδιότητα: α) $7\cdot 102$ \quad β) $8\cdot 99$ \quad γ) $25\cdot 12$ (σκέψου $12=10+2$).

\ex \textbf{Τα τσουβάλια.} Ο παππούς Γιάννης αγόρασε $6$ τσουβάλια λίπασμα των $25$ κιλών και μετά άλλα $4$ τσουβάλια των $25$ κιλών. Πόσα κιλά είναι όλα; Βρες το με δύο τρόπους.

\ex Υπολόγισε την τιμή: α) $40-5\cdot 6+8$ \quad β) $(40-5)\cdot 6+8$ \quad γ) $2\cdot(3+4)+5\cdot 2$ \quad δ) $3+2^3\cdot 5$.

\zoni{Το σκέφτομαι}{ερωτήσεις με λίγη περισσότερη σκέψη}

\ex Ποιος είναι ο μικρότερος τετραψήφιος και ποιος ο μεγαλύτερος τριψήφιος φυσικός; Πόσοι διψήφιοι αριθμοί υπάρχουν; Πόσοι από αυτούς είναι άρτιοι;

\ex Ο Βασίλης λέει: «Στην αφαίρεση ισχύει η αντιμεταθετική: $8-5=5-8$». Εξήγησε γιατί κάνει λάθος. Μετά δοκίμασε αν ισχύει η προσεταιριστική, δηλαδή αν $(8-3)-2=8-(3-2)$.

\zoni{Για το σπίτι}{μία ακόμη}

\ex Φτιάξε με τα ψηφία $5,0,8,2$ (κάθε ένα μία φορά) τον μεγαλύτερο και τον μικρότερο τετραψήφιο αριθμό. Πόση είναι η διαφορά τους;
